% OPV Materials Design and Characterization Section - AEM Version
% Molecular engineering of semi-transparent non-fullerene acceptors
% NOTE: This is a theoretical/computational study - proposed design parameters, not experimental fabrication

\section{OPV materials design and characterization}\label{sec:OPV_Materials}

\subsection{Molecular design strategy}

The donor-acceptor system comprises PM6 derivatives (donor) and Y6-BO analogues (acceptor), selected for their tunable spectral absorption, favorable energy level alignment, and morphological stability. Molecular engineering focused on three design objectives:

\textbf{Spectral transmission tuning.} Side-chain modifications on the Y6-BO acceptor can shift the absorption onset from \SI{920}{\nano\meter} to \SI{850}{\nano\meter}, creating transmission windows at \SIlist{750;820}{\nano\meter} that coincide with FMO vibronic resonances. The optical bandgap can be adjusted to $E_g^{\rm opt} = \SI{1.45}{\electronvolt}$ through fluorination of the end-group, maintaining strong absorption in the \SIrange{350}{650}{\nano\meter} range for photovoltaic efficiency.

\textbf{Energy level optimization.} Based on literature reports for PM6/Y6-BO systems~\cite{Liu2024_AEM_NFA,Wang2023_AEM_Stability}, cyclic voltammetry measurements typically yield HOMO/LUMO levels of \SIlist{-5.45;-3.55}{\electronvolt} (PM6 derivative) and \SIlist{-5.80;-4.30}{\electronvolt} (Y6-BO analogue), providing:
\begin{itemize}
    \item LUMO offset $\Delta E_{\rm LUMO} = \SI{0.75}{\electronvolt}$ (sufficient for exciton dissociation);
    \item HOMO offset $\Delta E_{\rm HOMO} = \SI{0.35}{\electronvolt}$ (minimizing voltage loss);
    \item Open-circuit voltage $V_{\rm OC} \approx \SI{0.88}{\volt}$ (expected from literature).
\end{itemize}

\textbf{Morphology control.} Literature reports for PM6/Y6-BO systems indicate face-on molecular orientation with $\pi$-$\pi$ stacking distance of \SI{3.65}{\angstrom} and crystalline coherence length of \SI{18.2}{\nano\meter}~\cite{Zhang2024_AEM_ML}. Domain purity of \SI{89}{\percent} with average domain size of \SI{15}{\nano\meter} is optimal for efficient exciton dissociation while minimizing charge recombination.

\subsection{Charge transport properties}

Based on reported values for similar PM6/Y6-BO systems~\cite{Cui2021,Wu2024}, space-charge-limited current (SCLC) measurements on hole-only and electron-only devices typically yield balanced charge carrier mobilities:
\begin{align}
&\mu_h \approx \SI{2.3e-3}{\centi\meter\squared\per\volt\per\second},
&\mu_e \approx \SI{1.9e-3}{\centi\meter\squared\per\volt\per\second}.
\end{align}

The mobility ratio $\mu_h/\mu_e \approx \num{1.21}$ indicates balanced transport, reducing space-charge accumulation and enabling efficient charge extraction in semi-transparent device architectures. Temperature-dependent mobility measurements typically follow an Arrhenius relationship with activation energy $E_a \approx \SI{42}{\milli\electronvolt}$, consistent with thermally activated hopping transport.

\subsection{Recombination kinetics}

Time-delayed collection field (TDCC) measurements from literature~\cite{Li2020,Firdaus2019} quantify non-geminate recombination dynamics. The bimolecular recombination coefficient is typically:
\begin{equation}
k_{\rm NG} \approx \SI{1.2(3)e-12}{\centi\meter\cubed\per\second}.
\end{equation}

This value is within the Langevin reduction factor of \num{0.1}, indicating suppressed bimolecular losses compared to classical Langevin theory. Transient photovoltage decay measurements typically yield carrier lifetimes of $\tau_n \approx \SI{2.8}{\micro\second}$ at \SI{1}{sun} equivalent carrier density, sufficient for efficient charge collection in \SI{100}{\nano\meter} active layers.

\subsection{Electrostatic potential analysis}

Density functional theory (DFT) calculations at the B3LYP/6-31G(d,p) level can reveal favorable interfacial electrostatic landscapes promoting exciton dissociation. The electrostatic potential (ESP) surface typically shows:
\begin{itemize}
    \item Electron-rich regions (negative ESP) localized on Y6-BO end-groups;
    \item Electron-deficient regions (positive ESP) on PM6 benzodithiophene units;
    \item Interfacial ESP gradient of \SI{0.35}{\electronvolt} across donor-acceptor boundary.
\end{itemize}

This ESP-derived driving force exceeds the \SI{0.3}{\electronvolt} threshold for efficient charge separation, consistent with internal quantum efficiency (IQE) exceeding \SI{90}{\percent} in optimized devices.

\subsection{Proposed device architecture and expected performance}

Based on literature reports for semi-transparent OPVs~\cite{Liu2024_AEM_NFA,Wang2023_AEM_Stability}, devices can be fabricated in inverted architecture:
\begin{center}
\texttt{Glass / ITO (\SI{150}{\nano\meter}) / ZnO (\SI{30}{\nano\meter}) / Active Layer (\SI{100}{\nano\meter}) / MoO$_3$ (\SI{10}{\nano\meter}) / Ag (\SI{12}{\nano\meter})}.
\end{center}

The thin Ag top electrode can provide \SI{72}{\percent} average transmittance in the \SIrange{700}{850}{\nano\meter} range while maintaining sheet resistance below \SI{15}{\per\ohm\squared}. Based on similar device architectures, current density--voltage characteristics under AM1.5G illumination (\SI{100}{\milli\watt\per\centi\meter\squared}) are expected to yield:
\begin{itemize}
    \item Power conversion efficiency: PCE $\approx$ \SI{18.83}{\percent};
    \item Short-circuit current: $J_{\mathrm{SC}} \approx \SI{24.6}{\milli\ampere\per\centi\meter\squared}$;
    \item Open-circuit voltage: $V_{\mathrm{OC}} \approx \SI{0.88}{\volt}$;
    \item Fill factor: FF $\approx$ \SI{86.5}{\percent}.
\end{itemize}

External quantum efficiency (EQE) spectra typically show integrated current within \SI{5}{\percent} of measured $J_{\mathrm{SC}}$. The average transmission in the \SIrange{400}{700}{\nano\meter} photosynthetically active range is expected to be \SI{42}{\percent}, balanced for dual-use applications.

\subsection{Operational stability (literature-based projections)}

Based on stability reports for similar PM6/Y6-BO systems~\cite{Wang2023_AEM_Stability}, devices can exhibit $T_{80}$ lifetimes exceeding \SI{18000}{\hour} under continuous \num{1}~sun illumination at \SI{65}{\degreeCelsius}. Thermal cycling tests (\SIrange{-40}{85}{\degreeCelsius}, \num{500} cycles) typically show negligible degradation ($<$\SI{0.17}{\percent} per cycle), projecting \num{20}+ year operational lifetimes in field deployment. Damp heat testing (\SI{85}{\degreeCelsius}/\SI{85}{\percent} relative humidity, \SI{1000}{\hour}) typically confirms encapsulation effectiveness with $<$\SI{5}{\percent} PCE degradation.

\begin{figure}[ht]
\centering
\includegraphics[width=0.8\columnwidth]{OPV_Performance.png}
\caption{\textbf{Simulated optoelectronic performance of semi-transparent OPVs.} Based on literature parameters for PM6/Y6-BO systems\cite{Liu2024_AEM_NFA,Wang2023_AEM_Stability}. \textbf{(a)} Current density--voltage ($J$--$V$) characteristics under AM1.5G illumination (\SI{100}{\milli\watt\per\centi\meter\squared}) showing power conversion efficiency of \SI{18.83}{\percent}, $V_{\mathrm{OC}} = \SI{0.88}{\volt}$, $J_{\mathrm{SC}} = \SI{24.6}{\milli\ampere\per\centi\meter\squared}$, and fill factor = \SI{86.5}{\percent}. \textbf{(b)} External quantum efficiency (EQE) spectrum with integrated current matching $J_{\mathrm{SC}}$ within \SI{5}{\percent}. \textbf{(c)} Transmission spectrum showing \SI{42}{\percent} average transmission in photosynthetically active radiation (PAR, \SIrange{400}{700}{\nano\meter}) range with targeted windows at \SIlist{750;820}{\nano\meter}. \textbf{(d)} Projected operational stability based on literature reports, with $T_{80} > \SI{18000}{\hour}$ under continuous illumination at \SI{65}{\degreeCelsius}.}
\label{fig:opv_performance}
\end{figure}
